% !TEX program = xelatex
\documentclass[aspectratio=169,11pt,handout]{beamer}
\usetheme{Madrid}
\usepackage{fontspec}
\setsansfont{DejaVu Sans}
\setmainfont{DejaVu Serif}
\usepackage{polyglossia}
\setdefaultlanguage{turkish}
\usepackage{amsmath,amssymb}
\usepackage{booktabs}
\usepackage{multicol}
\usepackage{xcolor}
\usepackage{tikz}

% ---------- Renkler ----------
\definecolor{anamavi}{RGB}{0,45,98}
\definecolor{cozumyesil}{RGB}{0,110,50}
\definecolor{uyarikirmizi}{RGB}{180,20,20}
\setbeamercolor{structure}{fg=anamavi}
\setbeamercolor{frametitle}{bg=anamavi,fg=white}
\setbeamercolor{title}{bg=anamavi,fg=white}
\setbeamercolor{palette primary}{bg=anamavi,fg=white}
\setbeamercolor{palette secondary}{bg=anamavi!85!black,fg=white}
\setbeamercolor{palette tertiary}{bg=anamavi!70!black,fg=white}
\setbeamercolor{block title}{bg=anamavi,fg=white}
\setbeamercolor{block body}{bg=anamavi!7}
\setbeamercolor{block title example}{bg=cozumyesil,fg=white}
\setbeamercolor{block body example}{bg=cozumyesil!8}
\setbeamercolor{block title alerted}{bg=uyarikirmizi,fg=white}
\setbeamercolor{block body alerted}{bg=uyarikirmizi!7}
\setbeamertemplate{navigation symbols}{}
\setbeamertemplate{blocks}[rounded][shadow=false]

% ---------- Soru / Çözüm kalıbı ----------
% \soru{başlık}{soru metni}{çözüm metni}: önce soru, bir tıklamayla çözüm
\newcommand{\soru}[3]{%
\begin{frame}{#1}
  \begin{block}{Soru}\large #2\end{block}
  \pause
  \begin{exampleblock}{Çözüm}\large #3\end{exampleblock}
\end{frame}}
\newcommand{\ab}[1]{\overline{#1}}

\title[MMAT105 -- Hafta 6]{MMAT105 Genel Matematik}
\subtitle{Hafta 6: Birinci Dereceden İki Bilinmeyenli Denklemler ve Denklem Sistemleri}
\author{}
\date{2026--2027 Güz Dönemi}

\begin{document}

% =====================================================
\begin{frame}
\titlepage
\end{frame}

% (iletişim çerçevesi kaldırıldı)

\begin{frame}{Geçen Hafta Hatırlatma}
\begin{columns}[t]
\column{0.58\textwidth}
\textbf{Çözünüz (Hafta 5):}
\begin{enumerate}
\item \(3(x-2)+5=2x+7\) ise \(x=?\)
\item \(\dfrac{x}{2}+\dfrac{x}{3}=10\) ise \(x=?\)
\item Bir sayının \(3\) katının \(4\) fazlası, aynı sayının \(5\) katının \(6\) eksiğine eşittir. Sayı kaçtır?
\end{enumerate}
\column{0.36\textwidth}
\pause
\textbf{Cevaplar:}
\begin{enumerate}
\item \(3x-1=2x+7\Rightarrow x=8\)
\item \(\dfrac{5x}{6}=10\Rightarrow x=12\)
\item \(3x+4=5x-6\Rightarrow x=5\)
\end{enumerate}
\end{columns}
\vspace{6pt}
\begin{block}{Bu hafta}
Bilinmeyen sayısı ikiye çıkıyor: \(x\) ve \(y\). Tek bir denklem yetmez, \textbf{iki denklem} gerekir.
\end{block}
\end{frame}

\begin{frame}{Bugün}
\begin{columns}[t]
\column{0.32\textwidth}
\begin{block}{Ders 1}
\begin{enumerate}
\item \(ax+by+c=0\) biçimi
\item Sıralı ikili, çözüm kontrolü
\item Denklem sistemi
\item Yok etme yöntemi
\end{enumerate}
\end{block}
\column{0.32\textwidth}
\begin{block}{Ders 2}
\begin{enumerate}\setcounter{enumi}{4}
\item Yerine koyma yöntemi
\item Kesirli / ondalıklı sistemler
\item Simetrik sistemler, \(x+y\), \(x\cdot y\)
\item Özel durumlar, parametre
\end{enumerate}
\end{block}
\column{0.32\textwidth}
\begin{block}{Ders 3}
\begin{enumerate}\setcounter{enumi}{8}
\item Sözel ifadeden sisteme
\item Para, bilet, puan
\item İki basamaklı sayılar
\item Bacak sayısı, yaş (giriş)
\end{enumerate}
\end{block}
\end{columns}
\end{frame}

% =====================================================
\section{Ders 1: İki Bilinmeyenli Denklem ve Yok Etme}
\begin{frame}
\centering\Large\color{anamavi}\textbf{Ders 1}\\[6pt]\normalsize İki Bilinmeyenli Denklem, Denklem Sistemi, Yok Etme Yöntemi
\end{frame}

\subsection{İki bilinmeyenli denklem}
\begin{frame}{Birinci Dereceden İki Bilinmeyenli Denklem}
\begin{block}{Tanım}
\(a,b,c\) gerçek sayılar, \(a\neq 0\) ve \(b\neq 0\) olmak üzere
\[ax+by+c=0\]
biçimindeki denkleme \textbf{birinci dereceden iki bilinmeyenli denklem} denir.
\begin{itemize}
\item Denklemi sağlayan her \((x,y)\) \textbf{sıralı ikilisi} bir çözümdür.
\item Tek bir denklemin \textbf{sonsuz sayıda} çözümü vardır: \(x\)'e bir değer verilir, \(y\) bulunur.
\end{itemize}
\end{block}
\begin{alertblock}{Sık yapılan hatalar}
\((x,y)\) sıralıdır: \((2,3)\neq(3,2)\). \ İlk bileşen her zaman \(x\), ikinci bileşen \(y\)'dir.
\end{alertblock}
\end{frame}

\soru{İki Bilinmeyenli Denklem -- Örnek 1}
{\(2x+y=7\) denkleminin üç çözümünü yazınız. \((3,2)\) bir çözüm müdür?}
{\(y=7-2x\): \quad \(x=0\Rightarrow y=7\), \ \(x=1\Rightarrow y=5\), \ \(x=2\Rightarrow y=3\)\\[3pt]
Çözümler: \((0,7),\ (1,5),\ (2,3),\ \dots\) \ (sonsuz tane)\\[3pt]
\((3,2)\): \(2\cdot 3+2=8\neq 7\) \(\Rightarrow\) çözüm \textbf{değildir}. \ (Oysa \((2,3)\) çözümdür!)}

\soru{Çözüm Kontrolü -- Örnek 2}
{\((3,-1)\), \ \((-2,4)\), \ \((1,2)\) ikililerinden hangileri \(3x+2y-7=0\) denklemini sağlar?}
{\((3,-1)\): \ \(9-2-7=0\) \ \checkmark\\[2pt]
\((-2,4)\): \ \(-6+8-7=-5\neq 0\) \ \(\times\)\\[2pt]
\((1,2)\): \ \(3+4-7=0\) \ \checkmark\\[4pt]
Cevap: \((3,-1)\) ve \((1,2)\).}

\soru{Çözüm Kontrolü -- Örnek 3}
{\((2,k)\) ikilisi \(5x-3y=1\) denklemini sağlıyorsa \(k\) kaçtır?}
{\(x=2,\ y=k\) yazılır:
\[5\cdot 2-3k=1\ \Rightarrow\ 10-3k=1\ \Rightarrow\ 3k=9\ \Rightarrow\ k=3\]}

\soru{Çözüm Kontrolü -- Örnek 4}
{\((a,\ a+1)\) ikilisi \(4x-y=8\) denkleminin bir çözümü ise \(a\) kaçtır?}
{\[4a-(a+1)=8\ \Rightarrow\ 3a-1=8\ \Rightarrow\ a=3\]
Çözüm ikilisi: \((3,4)\). \ Kontrol: \(12-4=8\).}

\soru{Doğal Sayı Çözümleri -- Örnek 5}
{\(2x+3y=24\) denklemini sağlayan kaç tane \((x,y)\) doğal sayı ikilisi vardır? \ (\(\mathbb{N}=\{0,1,2,\dots\}\))}
{\(2x=24-3y\) çift olmalı \(\Rightarrow\) \(y\) çift; \ \(3y\le 24\Rightarrow y\le 8\).\\[3pt]
\(y=0,2,4,6,8\ \Rightarrow\ x=12,9,6,3,0\)\\[3pt]
\((12,0),(9,2),(6,4),(3,6),(0,8)\) \(\Rightarrow\) \textbf{5 tane}.}

\subsection{Denklem sistemi ve yok etme}
\begin{frame}{Denklem Sistemi ve Yok Etme Yöntemi}
\begin{block}{Denklem sistemi}
\[\begin{cases}a_1x+b_1y=c_1\\ a_2x+b_2y=c_2\end{cases}\]
Sistemin çözümü: \textbf{iki denklemi birlikte} sağlayan \((x,y)\) ikilisidir.
\end{block}
\begin{block}{Yok etme yöntemi}
\begin{enumerate}
\item Bir bilinmeyenin katsayılarını (gerekirse denklemleri sabitlerle çarparak) eşit ya da zıt yap.
\item Denklemleri taraf tarafa topla (zıt ise) ya da çıkar (eşit ise): bir bilinmeyen yok olur.
\item Bulunan değeri denklemlerden birinde yerine koy; sonra iki denklemde de kontrol et.
\end{enumerate}
\end{block}
\end{frame}

\begin{frame}{Yok Etme -- Sık Yapılan Hatalar}
\begin{alertblock}{Sık yapılan hatalar}
\begin{itemize}
\item Denklemi bir sayıyla çarparken \textbf{sağ tarafı unutmak}: \(\ 2\cdot(x-y=3)\Rightarrow 2x-2y=\mathbf{6}\).
\item Çıkarırken işareti dağıtmamak: \((2x+3y)-(2x-y)=3y+y=4y\) \ (\(2y\) değil).
\item Yalnızca \(x\)'i bulup durmak: soru genellikle \(y\)'yi, \(x+y\)'yi ya da \(x\cdot y\)'yi sorar.
\item Kontrolü atlamak: bulunan ikili \textbf{iki denklemi de} sağlamalıdır.
\end{itemize}
\end{alertblock}
\end{frame}

\soru{Yok Etme -- Örnek 6}
{\(\begin{cases}x+y=12\\ x-y=4\end{cases}\) \quad sisteminin çözümü nedir?}
{\(y\)'lerin katsayıları zıt: taraf tarafa toplayalım.
\[2x=16\ \Rightarrow\ x=8,\qquad 8+y=12\ \Rightarrow\ y=4\]
Çözüm: \((8,4)\). \ Kontrol: \(8-4=4\) \checkmark}

\soru{Yok Etme -- Örnek 7}
{\(\begin{cases}2x+3y=13\\ 2x-y=1\end{cases}\) \quad sisteminin çözümü nedir?}
{\(x\)'lerin katsayıları eşit: taraf tarafa çıkaralım.
\[(2x+3y)-(2x-y)=13-1\ \Rightarrow\ 4y=12\ \Rightarrow\ y=3\]
\(2x-3=1\Rightarrow x=2\). \quad Çözüm: \((2,3)\).}

\soru{Yok Etme -- Örnek 8}
{\(\begin{cases}3x+2y=16\\ x-2y=0\end{cases}\) \quad ise \(x\) ve \(y\) kaçtır?}
{Topla: \(4x=16\Rightarrow x=4\).\\[3pt]
\(4-2y=0\Rightarrow y=2\). \quad Çözüm: \((4,2)\).\\[3pt]
Kontrol: \(3\cdot 4+2\cdot 2=16\) \checkmark}

\soru{Yok Etme -- Örnek 9 (Çarparak eşitleme)}
{\(\begin{cases}3x+4y=10\\ 2x-3y=1\end{cases}\) \quad sisteminin çözümü nedir?}
{\(y\)'yi yok edelim: 1. denklem \(\times 3\), 2. denklem \(\times 4\):
\[9x+12y=30,\qquad 8x-12y=4\]
Topla: \(17x=34\Rightarrow x=2\); \quad \(6+4y=10\Rightarrow y=1\). \quad Çözüm: \((2,1)\).}

\soru{Yok Etme -- Örnek 10}
{\(\begin{cases}5x+2y=4\\ 3x-4y=18\end{cases}\) \quad sisteminin çözümü nedir?}
{1. denklem \(\times 2\): \(10x+4y=8\). \ 2. denklemle topla:
\[13x=26\ \Rightarrow\ x=2\]
\(10+2y=4\Rightarrow y=-3\). \quad Çözüm: \((2,-3)\).\\[2pt]
Kontrol: \(6-4\cdot(-3)=18\) \checkmark}

\soru{Önce Düzenle -- Örnek 11}
{\(\begin{cases}2(x+1)-y=5\\ x+3(y-2)=-1\end{cases}\) \quad sisteminin çözümü nedir?}
{Düzenle: \(2x-y=3\) \ ve \ \(x+3y=5\).\\[3pt]
1. denklem \(\times 3\): \(6x-3y=9\); \ topla: \(7x=14\Rightarrow x=2\).\\[3pt]
\(4-y=3\Rightarrow y=1\). \quad Çözüm: \((2,1)\).}

\soru{Yok Etme -- Örnek 12}
{\(\begin{cases}4x-3y=-1\\ 6x+5y=27\end{cases}\) \quad ise \(x\cdot y\) kaçtır?}
{\(x\)'i yok edelim: 1. denklem \(\times 3\), 2. denklem \(\times 2\):
\[12x-9y=-3,\qquad 12x+10y=54\]
Çıkar: \(19y=57\Rightarrow y=3\); \quad \(4x-9=-1\Rightarrow x=2\).
\[x\cdot y=2\cdot 3=6\]}

\begin{frame}{Ders 1 -- Sınıfta Siz Çözün}
\begin{columns}[t]
\column{0.58\textwidth}
\begin{enumerate}
\item \((-1,m)\) ikilisi \(3x+2y=5\) denklemini sağlıyorsa \(m=?\)
\item \(x+y=9\), \ \(x-y=3\) ise \((x,y)=?\)
\item \(3x-y=7\), \ \(x+y=5\) ise \((x,y)=?\) \(\star\)
\item \(2x+5y=1\), \ \(3x-2y=11\) ise \((x,y)=?\) \(\star\)
\item \(x+2y=10\) denklemini sağlayan kaç \((x,y)\) doğal sayı ikilisi vardır? \(\star\)
\end{enumerate}
\column{0.36\textwidth}
\pause
\textbf{Cevaplar:}
\begin{enumerate}
\item \(m=4\)
\item \((6,3)\)
\item \((3,2)\)
\item \((3,-1)\)
\item \(6\) \ (\(y=0,1,\dots,5\))
\end{enumerate}
\end{columns}
\vspace{4pt}{\footnotesize \(\star\) işaretli soruların adım adım çözümü sonraki slaytlardadır.}
\end{frame}

% ---- Ders 1 alıştırma çözümleri ----
\soru{Ders 1 -- Alıştırma Çözümü (Soru 3)}
{\(\begin{cases}3x-y=7\\ x+y=5\end{cases}\) \quad sisteminin çözümü nedir?}
{\(y\)'lerin katsayıları zıt (\(-1\) ve \(+1\)): taraf tarafa toplayalım.
\[(3x-y)+(x+y)=7+5\ \Rightarrow\ 4x=12\ \Rightarrow\ x=3\]
\(3+y=5\Rightarrow y=2\). \quad Çözüm: \((3,2)\).\\[3pt]
Kontrol: \(3\cdot 3-2=7\) \checkmark, \ \(3+2=5\) \checkmark}

\soru{Ders 1 -- Alıştırma Çözümü (Soru 4)}
{\(\begin{cases}2x+5y=1\\ 3x-2y=11\end{cases}\) \quad sisteminin çözümü nedir?}
{\(y\)'yi yok edelim: 1. denklem \(\times 2\), 2. denklem \(\times 5\):
\[4x+10y=2,\qquad 15x-10y=55\]
Topla: \(19x=57\Rightarrow x=3\). \quad \(6+5y=1\Rightarrow 5y=-5\Rightarrow y=-1\).\\[3pt]
Çözüm: \((3,-1)\). \ Kontrol: \(9-2\cdot(-1)=11\) \checkmark}

\soru{Ders 1 -- Alıştırma Çözümü (Soru 5)}
{\(x+2y=10\) denklemini sağlayan kaç \((x,y)\) doğal sayı ikilisi vardır?}
{\(x=10-2y\ge 0\) olmalı \(\Rightarrow\) \(2y\le 10\Rightarrow y\le 5\).\\[3pt]
\(y=0,1,2,3,4,5\ \Rightarrow\ x=10,8,6,4,2,0\)\\[3pt]
\((10,0),(8,1),(6,2),(4,3),(2,4),(0,5)\) \(\Rightarrow\) \textbf{6 tane}.}

% =====================================================
\section{Ders 2: Yerine Koyma, Özel Durumlar}
\begin{frame}
\centering\Large\color{anamavi}\textbf{Ders 2}\\[6pt]\normalsize Yerine Koyma Yöntemi, Kısa Yollar ve Özel Durumlar
\end{frame}

\subsection{Yerine koyma}
\begin{frame}{Yerine Koyma Yöntemi}
\begin{block}{Adımlar}
\begin{enumerate}
\item Denklemlerden birinde bir bilinmeyeni yalnız bırak: ör. \(y=\dots\)
\item Bu ifadeyi \textbf{diğer} denklemde yerine koy: tek bilinmeyenli denklem elde edilir.
\item Çöz, sonra bulunan değeri geri koyarak ikinci bilinmeyeni bul.
\end{enumerate}
Bir bilinmeyenin katsayısı \(1\) veya \(-1\) ise bu yöntem hızlıdır.
\end{block}
\begin{alertblock}{Sık yapılan hatalar}
Yerine koyarken \textbf{parantez} kullanın: \(3x+2y\) içinde \(y=2x-1\) \(\Rightarrow\) \(3x+2(2x-1)\), \ \(3x+4x-1\) değil!\\
İfadeyi elde ettiğiniz denklemin kendisine geri koymayın (\(0=0\) çıkar, bilgi vermez).
\end{alertblock}
\end{frame}

\soru{Yerine Koyma -- Örnek 13}
{\(\begin{cases}y=2x-1\\ 3x+2y=12\end{cases}\) \quad sisteminin çözümü nedir?}
{\[3x+2(2x-1)=12\ \Rightarrow\ 7x-2=12\ \Rightarrow\ x=2\]
\(y=2\cdot 2-1=3\). \quad Çözüm: \((2,3)\).}

\soru{Yerine Koyma -- Örnek 14}
{\(\begin{cases}x-3y=1\\ 2x-5y=4\end{cases}\) \quad sisteminin çözümü nedir?}
{1. denklemden \(x=3y+1\):
\[2(3y+1)-5y=4\ \Rightarrow\ y+2=4\ \Rightarrow\ y=2\]
\(x=3\cdot 2+1=7\). \quad Çözüm: \((7,2)\).}

\soru{Yerine Koyma -- Örnek 15}
{\(\begin{cases}x=2y\\ 3x+y=21\end{cases}\) \quad ise \(x+y\) kaçtır?}
{\[3(2y)+y=21\ \Rightarrow\ 7y=21\ \Rightarrow\ y=3,\quad x=6\]
\[x+y=9\]}

\subsection{Kesirli ve ondalıklı katsayılar}
\begin{frame}{Kesirli ve Ondalıklı Katsayılar}
\begin{block}{Önce temizle}
\begin{itemize}
\item \textbf{Kesirli:} Her denklemi kendi paydalarının EKOK'u ile çarp.
\item \textbf{Ondalıklı:} Her denklemi \(10\), \(100\), \dots\ ile çarp.
\end{itemize}
Sonra yok etme ya da yerine koyma ile çöz.
\end{block}
\begin{alertblock}{Sık yapılan hatalar}
Çarparken \textbf{her terimi} (sağ taraf dahil) çarpın: \(\dfrac{x}{2}+\dfrac{y}{3}=4\ \xrightarrow{\ \times 6\ }\ 3x+2y=24\) \ (\(3x+2y=4\) değil).
\end{alertblock}
\end{frame}

\soru{Kesirli Katsayılar -- Örnek 16}
{\(\begin{cases}\dfrac{x}{2}+\dfrac{y}{3}=4\\[8pt] \dfrac{x}{4}-\dfrac{y}{6}=0\end{cases}\) \quad sisteminin çözümü nedir?}
{1. denklem \(\times 6\): \(3x+2y=24\); \quad 2. denklem \(\times 12\): \(3x-2y=0\).\\[3pt]
Topla: \(6x=24\Rightarrow x=4\); \quad \(12+2y=24\Rightarrow y=6\).\\[3pt]
Çözüm: \((4,6)\). \ Kontrol: \(2+2=4\), \ \(1-1=0\) \checkmark}

\soru{Ondalıklı Katsayılar -- Örnek 17}
{\(\begin{cases}0{,}2x+0{,}5y=1{,}9\\ 0{,}3x-0{,}1y=0{,}3\end{cases}\) \quad sisteminin çözümü nedir?}
{\(\times 10\): \ \(2x+5y=19\) \ ve \ \(3x-y=3\ \Rightarrow\ y=3x-3\).
\[2x+5(3x-3)=19\ \Rightarrow\ 17x=34\ \Rightarrow\ x=2\]
\(y=3\cdot 2-3=3\). \quad Çözüm: \((2,3)\).}

\subsection{Simetrik sistemler}
\begin{frame}{Kısa Yol: Simetrik Sistemler}
\begin{block}{Katsayılar yer değiştirmişse}
\[\begin{cases}ax+by=m\\ bx+ay=n\end{cases}\quad(a\neq b)\]
\begin{itemize}
\item Topla: \((a+b)(x+y)=m+n\ \Rightarrow\ x+y=\dfrac{m+n}{a+b}\) \quad (\(a+b\neq 0\))
\item Çıkar: \((a-b)(x-y)=m-n\ \Rightarrow\ x-y=\dfrac{m-n}{a-b}\)
\end{itemize}
Sonra \(x+y\) ve \(x-y\)'den \(x\) ve \(y\) bulunur: \(x=\dfrac{(x+y)+(x-y)}{2}\).
\end{block}
\begin{alertblock}{İpucu}
Soru \(x+y\), \(x-y\) ya da \(x\cdot y\) istiyorsa önce denklemleri toplamayı / çıkarmayı deneyin; çoğu zaman \(x\) ve \(y\)'yi tek tek bulmaya gerek kalmaz.
\end{alertblock}
\end{frame}

\soru{Simetrik Sistem -- Örnek 18}
{\(\begin{cases}3x+2y=12\\ 2x+3y=13\end{cases}\) \quad sisteminin çözümü nedir?}
{Topla: \(5x+5y=25\Rightarrow x+y=5\).\\[3pt]
Çıkar: \(x-y=-1\).\\[3pt]
\(2x=5+(-1)=4\Rightarrow x=2\), \ \(y=3\). \quad Çözüm: \((2,3)\).}

\soru{Simetrik Sistem -- Örnek 19}
{\(\begin{cases}7x+4y=37\\ 4x+7y=40\end{cases}\) \quad ise \(x\cdot y\) kaçtır?}
{Topla: \(11(x+y)=77\Rightarrow x+y=7\).\\[3pt]
Çıkar: \(3(x-y)=-3\Rightarrow x-y=-1\).\\[3pt]
\(x=3\), \ \(y=4\) \(\Rightarrow\) \(x\cdot y=12\).}

\soru{İfade İstenmesi -- Örnek 20}
{\(2x+y=11\) ve \(x+2y=7\) olduğuna göre \(x+y\) kaçtır?}
{\(x\) ve \(y\)'yi bulmadan, iki denklemi toplayalım:
\[3x+3y=18\ \Rightarrow\ x+y=6\]
(İstenirse: çıkarınca \(x-y=4\), buradan \(x=5,\ y=1\).)}

\soru{İfade İstenmesi -- Örnek 21}
{\(x+y=8\) ve \(x-y=2\) olduğuna göre \(x\cdot y\) ve \(\dfrac{x}{y}\) kaçtır?}
{Topla: \(2x=10\Rightarrow x=5\); \quad \(y=3\).
\[x\cdot y=15,\qquad \frac{x}{y}=\frac{5}{3}\]}

\subsection{Özel durumlar}
\begin{frame}{Özel Durumlar: Kaç Çözüm Var?}
\begin{block}{\(a_1x+b_1y=c_1\), \ \(a_2x+b_2y=c_2\) sistemi \quad (oranlar tanımlı)}
\[\renewcommand{\arraystretch}{1.2}
\begin{array}{lll}
\text{Koşul} & \text{Çözüm sayısı} & \text{Doğrular}\\\hline
\dfrac{a_1}{a_2}\neq\dfrac{b_1}{b_2} & \text{tek çözüm} & \text{kesişen}\\[10pt]
\dfrac{a_1}{a_2}=\dfrac{b_1}{b_2}\neq\dfrac{c_1}{c_2} & \text{çözüm yok} & \text{paralel}\\[10pt]
\dfrac{a_1}{a_2}=\dfrac{b_1}{b_2}=\dfrac{c_1}{c_2} & \text{sonsuz çözüm} & \text{çakışık}
\end{array}\]
\end{block}
\begin{alertblock}{Sık yapılan hatalar}
İki denklemi de \textbf{aynı biçimde} (\(ax+by=c\)) yazın; \(x\) katsayısını \(x\) katsayısıyla, \(y\)'yi \(y\)'yle, sabiti sabitle oranlayın. Tek çözümde \(\frac{c_1}{c_2}\)'ye bakmaya gerek yoktur.
\end{alertblock}
\end{frame}

\begin{frame}{Grafik Yorumu: Her Denklem Bir Doğrudur}
\begin{columns}[t]
\column{0.32\textwidth}
\centering
\begin{tikzpicture}[scale=0.55]
\draw[->,gray] (-0.5,0)--(5,0) node[below]{\scriptsize \(x\)};
\draw[->,gray] (0,-1.5)--(0,5) node[left]{\scriptsize \(y\)};
\draw[thick,anamavi] (-0.5,4.5)--(4.5,-0.5);
\draw[thick,uyarikirmizi] (1,-1)--(5,3);
\fill (3,1) circle (3pt);
\node[above right] at (3,1) {\scriptsize \((3,1)\)};
\end{tikzpicture}\\
\textbf{Kesişen}: tek çözüm\\
{\small \(x+y=4\), \ \(x-y=2\)}
\column{0.32\textwidth}
\centering
\begin{tikzpicture}[scale=0.55]
\draw[->,gray] (-0.5,0)--(5,0) node[below]{\scriptsize \(x\)};
\draw[->,gray] (0,-1.5)--(0,5) node[left]{\scriptsize \(y\)};
\draw[thick,anamavi] (-0.5,4.5)--(4.5,-0.5);
\draw[thick,uyarikirmizi] (-0.5,1.5)--(2.5,-1.5);
\end{tikzpicture}\\
\textbf{Paralel}: çözüm yok\\
{\small \(x+y=4\), \ \(x+y=1\)}
\column{0.32\textwidth}
\centering
\begin{tikzpicture}[scale=0.55]
\draw[->,gray] (-0.5,0)--(5,0) node[below]{\scriptsize \(x\)};
\draw[->,gray] (0,-1.5)--(0,5) node[left]{\scriptsize \(y\)};
\draw[line width=3pt,uyarikirmizi!50] (-0.5,4.5)--(4.5,-0.5);
\draw[thick,anamavi] (-0.5,4.5)--(4.5,-0.5);
\end{tikzpicture}\\
\textbf{Çakışık}: sonsuz çözüm\\
{\small \(x+y=4\), \ \(2x+2y=8\)}
\end{columns}
\vspace{6pt}
\begin{block}{Sezgi}
Bir denklemin çözümleri bir doğru oluşturur. Sistemin çözümü, iki doğrunun \textbf{ortak noktalarıdır}.
\end{block}
\end{frame}

\soru{Özel Durumlar -- Örnek 22}
{Aşağıdaki sistemlerin kaç çözümü vardır?\\
(a) \(2x-3y=5,\ 4x-6y=10\) \quad (b) \(2x-3y=5,\ 4x-6y=7\) \quad (c) \(2x-3y=5,\ x+y=5\)}
{(a) \(\frac{2}{4}=\frac{-3}{-6}=\frac{5}{10}=\frac12\) \(\Rightarrow\) \textbf{sonsuz} çözüm (çakışık)\\[3pt]
(b) \(\frac{2}{4}=\frac{-3}{-6}=\frac12\neq\frac{5}{7}\) \(\Rightarrow\) çözüm \textbf{yok} (paralel)\\[3pt]
(c) \(\frac{2}{1}\neq\frac{-3}{1}\) \(\Rightarrow\) \textbf{tek} çözüm: \((4,1)\)}

\soru{Parametreli Soru -- Örnek 23}
{\(\begin{cases}3x+my=5\\ 6x+4y=7\end{cases}\) \quad sisteminin çözümü yoksa \(m\) kaçtır?}
{Çözüm yok \(\Rightarrow\) \(\dfrac{3}{6}=\dfrac{m}{4}\neq\dfrac{5}{7}\).
\[\frac{m}{4}=\frac12\ \Rightarrow\ m=2\]
Kontrol: \(\frac12\neq\frac57\) \checkmark \ (Sistem: \(3x+2y=5\) ve \(3x+2y=3{,}5\): çelişki.)}

\soru{Parametreli Soru -- Örnek 24}
{\(\begin{cases}2x-y=3\\ ax+3y=b\end{cases}\) \quad sisteminin sonsuz çözümü varsa \(a+b\) kaçtır?}
{\[\frac{2}{a}=\frac{-1}{3}=\frac{3}{b}\]
\(\frac{2}{a}=-\frac13\Rightarrow a=-6\); \quad \(\frac{3}{b}=-\frac13\Rightarrow b=-9\).
\[a+b=-15\]
(2. denklem, 1. denklemin \(-3\) katıdır: \(-6x+3y=-9\).)}

\soru{Parametreli Soru -- Örnek 25}
{\(\begin{cases}(m-1)x+2y=4\\ 3x+y=5\end{cases}\) \quad sisteminin tek çözümü olması için \(m\) hangi değeri \textbf{alamaz}?}
{Tek çözüm \(\Rightarrow\) \(\dfrac{m-1}{3}\neq\dfrac{2}{1}\ \Rightarrow\ m-1\neq 6\ \Rightarrow\ m\neq 7\).\\[3pt]
\(m=7\) için: \(6x+2y=4\) ve \(6x+2y=10\) \(\Rightarrow\) çözüm yok.\\[3pt]
Cevap: \(m=7\) olamaz.}

\soru{Parametreli Soru -- Örnek 26}
{\(\begin{cases}x+2y=5\\ 3x+ky=15\end{cases}\) \quad sisteminin sonsuz çözümü varsa \(k\) kaçtır? Bu durumda \((1,2)\) bir çözüm müdür?}
{\(\dfrac13=\dfrac{2}{k}=\dfrac{5}{15}\) \(\Rightarrow\) \(k=6\).\\[3pt]
Denklemler aynı doğruyu gösterir; \((1,2)\): \(1+4=5\) \checkmark \ ve \(3+12=15\) \checkmark\\[3pt]
Evet, \((1,2)\) çözümdür (sonsuz çözümden biri).}

\begin{frame}{Ders 2 -- Sınıfta Siz Çözün}
\begin{columns}[t]
\column{0.58\textwidth}
\begin{enumerate}
\item \(y=x+2\), \ \(2x+y=11\) ise \((x,y)=?\)
\item \(5x+3y=21\), \ \(3x+5y=19\) ise \(x+y=?\), \ \(x-y=?\) \(\star\)
\item \(\dfrac{x}{3}+\dfrac{y}{2}=3\), \ \(x-y=-1\) ise \((x,y)=?\) \(\star\)
\item \(4x+ky=3\), \ \(2x+3y=5\) sisteminin çözümü yoksa \(k=?\) \(\star\)
\item \(x-2y=1\), \ \(3x-6y=3\) sisteminin kaç çözümü vardır?
\end{enumerate}
\column{0.36\textwidth}
\pause
\textbf{Cevaplar:}
\begin{enumerate}
\item \((3,5)\)
\item \(5\); \ \(1\)
\item \((3,4)\)
\item \(k=6\)
\item sonsuz çözüm
\end{enumerate}
\end{columns}
\vspace{4pt}{\footnotesize \(\star\) işaretli soruların adım adım çözümü sonraki slaytlardadır.}
\end{frame}

% ---- Ders 2 alıştırma çözümleri ----
\soru{Ders 2 -- Alıştırma Çözümü (Soru 2)}
{\(5x+3y=21\) ve \(3x+5y=19\) ise \(x+y\) ve \(x-y\) kaçtır?}
{Simetrik sistem. Topla: \(8x+8y=40\Rightarrow x+y=5\).\\[3pt]
Çıkar: \(2x-2y=2\Rightarrow x-y=1\).\\[3pt]
(İstenirse: \(2x=6\Rightarrow x=3\), \(y=2\). Kontrol: \(15+6=21\) \checkmark)}

\soru{Ders 2 -- Alıştırma Çözümü (Soru 3)}
{\(\dfrac{x}{3}+\dfrac{y}{2}=3\) \ ve \ \(x-y=-1\) ise \((x,y)\) nedir?}
{1. denklem \(\times 6\): \(2x+3y=18\). \quad 2. denklemden \(x=y-1\):
\[2(y-1)+3y=18\ \Rightarrow\ 5y=20\ \Rightarrow\ y=4\]
\(x=4-1=3\). \quad Çözüm: \((3,4)\). \ Kontrol: \(1+2=3\) \checkmark}

\soru{Ders 2 -- Alıştırma Çözümü (Soru 4)}
{\(\begin{cases}4x+ky=3\\ 2x+3y=5\end{cases}\) \quad sisteminin çözümü yoksa \(k\) kaçtır?}
{Çözüm yok \(\Rightarrow\) \(\dfrac{4}{2}=\dfrac{k}{3}\neq\dfrac{3}{5}\).
\[\frac{k}{3}=2\ \Rightarrow\ k=6\]
Kontrol: \(2\neq\frac35\) \checkmark \ (2. denklem \(\times 2\): \(4x+6y=10\); oysa \(4x+6y=3\): çelişki.)}

% =====================================================
\section{Ders 3: Sözel Problemler}
\begin{frame}
\centering\Large\color{anamavi}\textbf{Ders 3}\\[6pt]\normalsize Sözel Problemlerden Denklem Sistemi Kurma
\end{frame}

\begin{frame}{Sözel Problemden Sisteme}
\begin{block}{Adımlar}
\begin{enumerate}
\item İstenen iki büyüklüğe harf ver: \(x=\dots\), \ \(y=\dots\) \ (ne olduklarını yaz!)
\item Metindeki \textbf{iki ayrı bilgiyi} iki denkleme çevir.
\item Sistemi çöz; cevabın \textbf{sorulan şey} olduğundan emin ol.
\item Sonucu metinde kontrol et (sayılar tam sayı mı, rakam mı, anlamlı mı?).
\end{enumerate}
\end{block}
\begin{block}{Sık kullanılan kalıplar}
\begin{itemize}
\item Adet ve toplam tutar: \(x+y=\text{adet}\), \ \(p\,x+q\,y=\text{tutar}\)
\item İki basamaklı sayı: \(\ab{ab}=10a+b\), \ \(\ab{ab}-\ab{ba}=9(a-b)\), \ \(\ab{ab}+\ab{ba}=11(a+b)\)
\end{itemize}
\end{block}
\end{frame}

\soru{İki Sayı -- Örnek 27}
{İki sayının toplamı \(50\), farkı \(14\)'tür. Büyük sayı kaçtır?}
{Büyük \(x\), küçük \(y\): \quad \(x+y=50\), \ \(x-y=14\).
\[2x=64\ \Rightarrow\ x=32,\qquad y=18\]
Cevap: \(32\).}

\soru{İki Sayı -- Örnek 28}
{İki sayıdan büyüğü, küçüğünün \(3\) katından \(4\) fazladır. Sayıların toplamı \(44\) olduğuna göre büyük sayı kaçtır?}
{\(x=3y+4\) \ ve \ \(x+y=44\):
\[3y+4+y=44\ \Rightarrow\ 4y=40\ \Rightarrow\ y=10\]
\(x=34\). \quad Kontrol: \(34=3\cdot 10+4\) \checkmark}

\soru{Para -- Örnek 29}
{Bir cüzdanda yalnızca \(5\) TL'lik ve \(10\) TL'lik banknotlardan toplam \(17\) banknot ve \(120\) TL vardır. Kaç tane \(10\) TL'lik banknot vardır?}
{\(x\): 5 TL'lik, \ \(y\): 10 TL'lik sayısı.
\[x+y=17,\qquad 5x+10y=120\ \Rightarrow\ x+2y=24\]
Çıkar: \(y=7\), \ \(x=10\). \ Kontrol: \(50+70=120\) \checkmark\\
Cevap: \(7\).}

\soru{Bilet -- Örnek 30}
{Bir tiyatroda öğrenci bileti \(40\) TL, tam bilet \(70\) TL'dir. Satılan \(30\) bilet için \(1560\) TL toplanmıştır. Kaç öğrenci bileti satılmıştır?}
{\(x\): öğrenci, \ \(y\): tam. \quad \(x+y=30\), \ \(40x+70y=1560\Rightarrow 4x+7y=156\).
\[4(30-y)+7y=156\ \Rightarrow\ 3y=36\ \Rightarrow\ y=12\]
\(x=18\). \ Kontrol: \(720+840=1560\) \checkmark \quad Cevap: \(18\).}

\soru{Alışveriş -- Örnek 31}
{\(2\) kalem ile \(3\) defter \(85\) TL, \(3\) kalem ile \(2\) defter \(65\) TL'dir. \(1\) kalem ile \(1\) defter kaç TL'dir? Bir defter kaç TL'dir?}
{\(k\): kalem, \ \(d\): defter. \quad \(2k+3d=85\), \ \(3k+2d=65\).\\[3pt]
Topla: \(5(k+d)=150\Rightarrow k+d=30\) TL.\\[3pt]
Çıkar: \(d-k=20\) \(\Rightarrow\) \(d=25\) TL, \ \(k=5\) TL.}

\soru{Bacak Sayısı -- Örnek 32}
{Bir çiftlikteki tavuk ve tavşanların toplam \(20\) başı ve \(56\) bacağı vardır. Kaç tavşan vardır?}
{\(t\): tavuk, \ \(r\): tavşan. \quad \(t+r=20\), \ \(2t+4r=56\Rightarrow t+2r=28\).
\[\text{Çıkar: } r=8,\qquad t=12\]
Pratik yol: Hepsi tavuk olsaydı \(40\) bacak olurdu; fazla \(16\) bacak, her tavşan \(+2\) \(\Rightarrow\) \(8\) tavşan.}

\soru{Puan -- Örnek 33}
{\(20\) soruluk bir sınavda doğru cevap \(+5\), yanlış cevap \(-2\) puandır. Bütün soruları cevaplayan Ece \(72\) puan almıştır. Ece kaç soruyu doğru cevaplamıştır?}
{\(d\): doğru, \ \(y\): yanlış. \quad \(d+y=20\), \ \(5d-2y=72\).
\[5d-2(20-d)=72\ \Rightarrow\ 7d=112\ \Rightarrow\ d=16\]
\(y=4\). \ Kontrol: \(80-8=72\) \checkmark}

\soru{İki Basamaklı Sayı -- Örnek 34}
{İki basamaklı \(\ab{ab}\) sayısının rakamları toplamı \(11\)'dir. Rakamları yer değiştirince sayı \(27\) artıyor. Bu sayı kaçtır?}
{\(a+b=11\) \ ve \ \(\ab{ba}-\ab{ab}=9(b-a)=27\Rightarrow b-a=3\).
\[2b=14\ \Rightarrow\ b=7,\quad a=4\]
Sayı: \(47\). \ Kontrol: \(74-47=27\) \checkmark}

\soru{İki Basamaklı Sayı -- Örnek 35}
{İki basamaklı \(\ab{ab}\) sayısı, rakamları toplamının \(7\) katına eşittir. Onlar basamağındaki rakam birler basamağındakinden \(3\) fazladır. Sayı kaçtır?}
{\(10a+b=7(a+b)\Rightarrow 3a=6b\Rightarrow a=2b\).\\[3pt]
\(a-b=3\Rightarrow 2b-b=3\Rightarrow b=3,\ a=6\).\\[3pt]
Sayı: \(63\). \ Kontrol: \(7\cdot(6+3)=63\) \checkmark}

\soru{Yaş (Giriş) -- Örnek 36}
{Ali ile babasının bugünkü yaşları toplamı \(50\)'dir. Babası Ali'den \(26\) yaş büyüktür. Ali kaç yaşındadır? \(5\) yıl sonra yaşları toplamı kaç olur?}
{\(a\): Ali, \ \(b\): baba. \quad \(a+b=50\), \ \(b-a=26\).
\[2b=76\Rightarrow b=38,\qquad a=12\]
\(5\) yıl sonra her biri \(5\) yaş büyür: \(50+2\cdot 5=60\).\\
{\normalsize (Yaş problemleri 13. haftada ayrıntılı.)}}

\begin{frame}{Ders 3 -- Sınıfta Siz Çözün}
\begin{columns}[t]
\column{0.62\textwidth}
\begin{enumerate}
\item İki sayının toplamı \(40\), büyüğü küçüğünün \(4\) katıdır. Küçük sayı?
\item \(20\) TL ve \(50\) TL'lik toplam \(12\) banknot \(390\) TL ediyor. Kaç tane \(50\) TL'lik var? \(\star\)
\item Tavuk ve koyunların toplam \(30\) başı, \(84\) bacağı var. Kaç koyun var? \(\star\)
\item \(\ab{ab}\): \(a+b=9\), \ \(\ab{ab}-\ab{ba}=45\). \ \(\ab{ab}=?\) \(\star\)
\item \(3\) simit \(+\) \(2\) çay \(55\) TL, \(1\) simit \(+\) \(2\) çay \(25\) TL. Bir çay kaç TL?
\end{enumerate}
\column{0.32\textwidth}
\pause
\textbf{Cevaplar:}
\begin{enumerate}
\item \(8\)
\item \(5\)
\item \(12\)
\item \(72\)
\item \(5\) TL
\end{enumerate}
\end{columns}
\vspace{4pt}{\footnotesize \(\star\) işaretli soruların adım adım çözümü sonraki slaytlardadır.}
\end{frame}

% ---- Ders 3 alıştırma çözümleri ----
\soru{Ders 3 -- Alıştırma Çözümü (Soru 2)}
{\(20\) TL'lik ve \(50\) TL'lik toplam \(12\) banknot \(390\) TL ediyor. Kaç tane \(50\) TL'lik banknot var?}
{\(x\): 20 TL'lik, \ \(y\): 50 TL'lik sayısı.
\[x+y=12,\qquad 20x+50y=390\ \Rightarrow\ 2x+5y=39\]
\(2(12-y)+5y=39\Rightarrow 3y=15\Rightarrow y=5\), \ \(x=7\).\\[2pt]
Kontrol: \(140+250=390\) \checkmark \quad Cevap: \(5\).}

\soru{Ders 3 -- Alıştırma Çözümü (Soru 3)}
{Tavuk ve koyunların toplam \(30\) başı ve \(84\) bacağı vardır. Kaç koyun vardır?}
{\(t\): tavuk, \ \(k\): koyun. \quad \(t+k=30\), \ \(2t+4k=84\Rightarrow t+2k=42\).
\[\text{Çıkar: } k=12,\qquad t=18\]
Kontrol: \(36+48=84\) \checkmark \quad Cevap: \(12\) koyun.}

\soru{Ders 3 -- Alıştırma Çözümü (Soru 4)}
{İki basamaklı \(\ab{ab}\) sayısında \(a+b=9\) ve \(\ab{ab}-\ab{ba}=45\) ise \(\ab{ab}\) kaçtır?}
{\(\ab{ab}-\ab{ba}=9(a-b)=45\Rightarrow a-b=5\).
\[a+b=9,\ a-b=5\ \Rightarrow\ 2a=14\ \Rightarrow\ a=7,\ b=2\]
\(\ab{ab}=72\). \ Kontrol: \(72-27=45\) \checkmark}

% =====================================================
\begin{frame}{Özet}
\begin{columns}[t]
\column{0.32\textwidth}
\begin{block}{Ders 1}
\begin{itemize}
\item \(ax+by+c=0\): sonsuz çözüm
\item Çözüm \(=\) sıralı ikili \((x,y)\)
\item Yok etme: katsayıyı eşitle, topla/çıkar
\end{itemize}
\end{block}
\column{0.32\textwidth}
\begin{block}{Ders 2}
\begin{itemize}
\item Yerine koyma: parantezle!
\item Kesir/ondalık: önce temizle
\item Simetrik: topla \(\to x+y\), çıkar \(\to x-y\)
\item Oranlar: tek / yok / sonsuz
\end{itemize}
\end{block}
\column{0.32\textwidth}
\begin{block}{Ders 3}
\begin{itemize}
\item İki bilgi \(\to\) iki denklem
\item Adet \(+\) tutar kalıbı
\item \(\ab{ab}\pm\ab{ba}\) kalıbı
\item Cevabı metinde kontrol et
\end{itemize}
\end{block}
\end{columns}
\vspace{6pt}
\begin{alertblock}{Gelecek hafta}
Basit eşitsizlikler (çözüm kümesi, sayı doğrusu) ve mutlak değer
\end{alertblock}
\end{frame}

\begin{frame}
\centering\Large İyi çalışmalar
\end{frame}

\end{document}
